%-------------------------------------------------------------------------------
% Mark J. Madler -- Resume (Fall 2026)
%-------------------------------------------------------------------------------
\documentclass[letterpaper,10pt]{article}

\usepackage[utf8]{inputenc}
\usepackage[T1]{fontenc}
\usepackage{lmodern}
\usepackage[hmargin=0.55in, vmargin=0.45in]{geometry}
\usepackage{titlesec}
\usepackage{enumitem}
\usepackage{tabularx}
\usepackage{xcolor}
\usepackage[hidelinks]{hyperref}

\definecolor{accent}{RGB}{0,51,102}

\pagestyle{empty}
\setlength{\parindent}{0pt}
\setlength{\tabcolsep}{0pt}
\raggedright

% Section headings: accent-colored small caps with a rule underneath
\titleformat{\section}{\color{accent}\large\scshape}{}{0pt}{}%
  [\vspace{-7pt}{\color{accent}\rule{\textwidth}{0.6pt}}]
\titlespacing*{\section}{0pt}{5pt}{3pt}

% Tight bullet lists
\setlist[itemize]{leftmargin=1.4em, itemsep=0pt, topsep=0pt, parsep=0pt,
                  label=\textbullet}

% \entry{Left bold}{Right}{Left italic}{Right italic}
\newcommand{\entry}[4]{%
  \begin{tabularx}{\textwidth}{X r}
    \textbf{#1} & #2 \\
    \textit{#3} & \textit{#4} \\
  \end{tabularx}\par\vspace{1pt}%
}

% \pub{Title}{Venue / status}{Authors}
\newcommand{\pub}[3]{%
  \textbf{#1}\par
  \begin{tabularx}{\textwidth}{X r}
    {\small #3} & \textit{#2} \\
  \end{tabularx}\par\vspace{1pt}%
}

% Bold my name in author lists
\newcommand{\me}{\textbf{M.~Madler}}

\begin{document}

%----------------------------------- HEADER ------------------------------------
\begin{center}
  {\Huge\scshape Mark J. Madler} \\[3pt]
  Boulder, CO \;$\cdot$\;
  (928) 925-4570 \;$\cdot$\;
  \href{mailto:madlermark@gmail.com}{madlermark@gmail.com} \;$\cdot$\;
  \href{https://github.com/MadlerMark}{github.com/MadlerMark}
\end{center}
\vspace{-2pt}

Personable and driven Computer Engineering PhD student researching programming
models, memory consistency, and verification for RDMA and CXL. Passionate about
teaching and engineering design.

%---------------------------------- EDUCATION ----------------------------------
\section{Education}

\begin{tabularx}{\textwidth}{X r}
  \textbf{University of Colorado Boulder} & Boulder, CO \\
  \textit{Ph.D., Computer Engineering \textendash{} Advisor: Dr.\ Joseph Izraelevitz}
    & \textit{Fall 2024 \textendash{} Present} \\
  \textit{M.S., Computer Engineering} & \textit{GPA: 3.97} \\
\end{tabularx}\par\vspace{1pt}

\entry{Gonzaga University}{Spokane, WA}
      {B.S., Computer Engineering}{GPA: 3.66}

%---------------------------------- EXPERIENCE ---------------------------------
\section{Experience}

\entry{University of Colorado Boulder}{August 2024 \textendash{} Present}
      {Graduate Research Assistant, CU Architecture Lab}{Boulder, CO}
\begin{itemize}
  \item Research RDMA programming models: encapsulation, memory consistency,
        and formally verified composable libraries.
\end{itemize}

\entry{Pacific Northwest National Laboratory}{May 2025 \textendash{} March 2026}
      {Research Intern \textendash{} Mentors: Bruce Palmer, Sayan Ghosh}{}
\begin{itemize}
  \item Compared one-sided, two-sided, and collective communication for
        distributed SpMM on NERSC Perlmutter.
  \item Studied the effects of communication models over CXL interconnects.
\end{itemize}

\entry{Kaiser Aluminum}{May \textendash{} December 2023}
      {Automation Engineering Intern}{}
\begin{itemize}
  \item Implemented industrial automation systems using Ignition SCADA.
  \item Diagnosed and fixed a process issue causing metal defects using
        Allen-Bradley PLCs.
  \item Managed network solutions on a large-scale OT network using vSphere.
\end{itemize}

\entry{Canyon AeroConnect}{May \textendash{} August 2022}
      {Engineering Intern}{}
\begin{itemize}
  \item Worked on a team to develop a tool for customizing a radio's
        functional options.
  \item Developed an internal tool for capturing and filtering telemetry data
        from radio systems.
\end{itemize}

%---------------------------------- TEACHING -----------------------------------
\section{Teaching}

\entry{ECEN 4313/5313: Concurrent Programming}{Fall 2026}
      {Teaching Assistant, University of Colorado Boulder}{}
\begin{itemize}
  \item Delivered several lectures on parallel programming concepts (coherence, OpenMP) and tools related to multithreaded debugging (perf).
\end{itemize}

%--------------------------------- PUBLICATIONS --------------------------------
\section{Publications}

\pub{A Verified High-Performance Composable Object Library for Remote Direct
     Memory Access}{POPL 2026}
    {G.~Ambal, G.~Hodgkins, \me, G.~Chockler, B.~Dongol, J.~Izraelevitz,
     A.~Raad, V.~Vafeiadis}

\pub{LOCURA: Lazy Consistency for RDMA Distributed Shared Memory}
    {EuroSys 2027 (in prep.)}
    {\me, G.~Ambal, E.~Karabacak, B.~Palmer, S.~Ghosh, J.~Izraelevitz}
\begin{itemize}
  \item RDMA DSM using lazy release consistency and multi-writer cachelines;
        formally proven correct.
  \item 1.9$\times$ faster than MPI RMA on Graph500; 3\textendash139$\times$
        faster than Argo DSM and 1.2\textendash6$\times$ faster than UPC++.
\end{itemize}

\pub{Understanding Performance Impacts of Communication Models on Sparse Matrix
     Multiplication}{IEEE CAL (in prep.)}
    {\me, B.~Palmer, S.~Ghosh, J.~Izraelevitz}
\begin{itemize}
  \item Shows one-sided communication wins by decoupling inter-process
        dependencies, not by curing load imbalance.
\end{itemize}

\pub{FLAME: Function-Level Analysis and Merging with Embeddings}
    {CGO 2027 (under review)}
    {H.~AlGhanim, \me, H.~Malahieude, J.~Izraelevitz}
\begin{itemize}
  \item Embedding-guided LLVM function merging: 11.02\% size reduction over
        \texttt{-Oz} vs.\ 9.06\% for F3M.
\end{itemize}

%----------------------------------- SKILLS ------------------------------------
\section{Skills}

\begin{tabularx}{\textwidth}{@{}l@{\hspace{1em}}X@{}}
  \textbf{Languages:}       & C, C++, Python, C\#, Assembly, MATLAB,
                              VHDL, Verilog \\
  \textbf{HPC \& Systems:}  & RDMA, MPI RMA, UPC++, Global Arrays, CXL,
                              pthreads, LLVM, eBPF, gem5 \\
  \textbf{Tools:}           & Linux, Git/GitHub, Docker, Make/CMake, \LaTeX,
                              Slurm, vSphere, Ignition, SolidWorks \\
\end{tabularx}

\end{document}
